\begin{proof}[Proof of \cref{obj:lem:clip-bias}]
Fix a bounded measurable covariate test function \(\varphi\), and write \(\bar e_q(x)=e_q(x;\bar\eta)\) for the clipped plug-in propensity of \cref{obj:def:clipped-propensity}, whose clipping level is fixed in \(q\in(0,1/2]\).  Since \(q\le 1/2\), we have \(q\le 1-q\), and the definition \(\bar e_q(x)=\min\{1-q,\max\{q,\bar e(x)\}\}\) therefore gives the two-sided bound
\[
q\le \bar e_q(x)\le 1-q
\qquad\text{for every }x,
\]
so both denominators are bounded away from zero: \(\bar e_q(x)\ge q>0\) and \(1-\bar e_q(x)\ge q>0\).  Together with \cref{obj:ass:bounded-outcome} (which gives \(|Y|\le1\) \(P\)-a.s.\ and \(|\mu_a|\le1\)) and the fixed plug-in regularity for this drift calculation (measurability of \(\bar e\), measurability of \(\bar\mu_0,\bar\mu_1\), and uniform boundedness of \(\bar\mu_0,\bar\mu_1\)), every integrand appearing below is measurable and bounded, hence \(P\)- and \(P_X\)-integrable; this is the only role of the boundedness hypotheses. % lean: clippedPropensity_pos % lean: one_sub_clippedPropensity_pos % lean: bounded_clipBias

\emph{Step 1: expand the score and pass to covariate integrals.}  Multiplying the score of \cref{obj:def:clipped-aipw-score} by \(\varphi(X)\) and splitting it into five terms,
\[
\varphi(X)\Gamma_q(O;\bar\eta)
=
\varphi(X)\{\bar\mu_1(X)-\bar\mu_0(X)\}
+\frac{AY\varphi(X)}{\bar e_q(X)}
-\frac{A\bar\mu_1(X)\varphi(X)}{\bar e_q(X)}
-\frac{(1-A)Y\varphi(X)}{1-\bar e_q(X)}
+\frac{(1-A)\bar\mu_0(X)\varphi(X)}{1-\bar e_q(X)} .
\]
% lean: clippedAIPWScore_test_expand_pointwise
Each of the five terms is integrable, so the integral of the sum splits into the sum of the integrals. % lean: hscore_split
The defining conditional-mean relations of the observed law, namely \(e_P(x)=P(A=1\mid X=x)\) and \(\mu_a(x)=E_P[Y\mid A=a,X=x]\), state that for every bounded measurable \(\psi\),
\[
E_P[\psi(X)]=\int \psi\,dP_X,
\qquad
E_P[A\,\psi(X)]=\int e_P\psi\,dP_X,
\qquad
E_P[(1-A)\,\psi(X)]=\int(1-e_P)\psi\,dP_X,
\]
\[
E_P[A\,Y\,\psi(X)]=\int e_P\mu_1\psi\,dP_X,
\qquad
E_P[(1-A)\,Y\,\psi(X)]=\int(1-e_P)\mu_0\psi\,dP_X .
\]
% lean: hI0 % lean: hI1 % lean: hI2 % lean: hI3 % lean: hI4
Applying these with \(\psi=\varphi/\bar e_q\), \(\psi=\bar\mu_1\varphi/\bar e_q\), \(\psi=\varphi/(1-\bar e_q)\), and \(\psi=\bar\mu_0\varphi/(1-\bar e_q)\) -- all bounded and measurable by the previous paragraph -- converts the five terms into \(P_X\)-integrals:
\[
\begin{aligned}
\int \varphi(X)\Gamma_q(O;\bar\eta)\,dP
&=\int \varphi(\bar\mu_1-\bar\mu_0)\,dP_X  \\
&\quad+\int e_P\mu_1\,\frac{\varphi}{\bar e_q}\,dP_X
      -\int e_P\bar\mu_1\,\frac{\varphi}{\bar e_q}\,dP_X \\
&\quad-\int (1-e_P)\mu_0\,\frac{\varphi}{1-\bar e_q}\,dP_X
      +\int (1-e_P)\bar\mu_0\,\frac{\varphi}{1-\bar e_q}\,dP_X .
\end{aligned}
\]
% lean: hscore_px

\emph{Step 2: collect the integrand pointwise.}  Fix \(x\).  Using \(\bar\mu_a=\mu_a+\Delta_a\) and \(\tau_P=\mu_1-\mu_0\), the two treated terms combine as
\[
\Delta_1(x)-\frac{e_P(x)\Delta_1(x)}{\bar e_q(x)}
=
\frac{\{\bar e_q(x)-e_P(x)\}\Delta_1(x)}{\bar e_q(x)},
\]
and the two control terms combine as
\[
-\Delta_0(x)+\frac{\{1-e_P(x)\}\Delta_0(x)}{1-\bar e_q(x)}
=
\frac{\{\bar e_q(x)-e_P(x)\}\Delta_0(x)}{1-\bar e_q(x)} .
\]
Both rearrangements are legitimate because \(\bar e_q(x)\) and \(1-\bar e_q(x)\) are nonzero.  Adding them to \(\tau_P(x)\) gives the pointwise identity
\[
\varphi(x)(\bar\mu_1-\bar\mu_0)(x)
+\frac{e_P\mu_1\varphi}{\bar e_q}(x)
-\frac{e_P\bar\mu_1\varphi}{\bar e_q}(x)
-\frac{(1-e_P)\mu_0\varphi}{1-\bar e_q}(x)
+\frac{(1-e_P)\bar\mu_0\varphi}{1-\bar e_q}(x)
=
\varphi(x)\{\tau_P(x)+b_q(x)\},
\]
with \(b_q\) exactly the function displayed in the statement. % lean: clipBias_integrand_collect_pointwise

\emph{Step 3: integrate.}  Both \(\varphi\tau_P\) and \(\varphi b_q\) are bounded measurable, hence \(P_X\)-integrable, so the integral of the right-hand side splits and Step 1 becomes
\[
\int \varphi(X)\Gamma_q(O;\bar\eta)\,dP
-
\int \varphi(x)\tau_P(x)\,dP_X(x)
=
\int \varphi(x)\,
\{\bar e_q(x)-e_P(x)\}
\left\{
\frac{\Delta_1(x)}{\bar e_q(x)}
+
\frac{\Delta_0(x)}{1-\bar e_q(x)}
\right\}\,dP_X(x)
=
\int \varphi\,b_q\,dP_X .
\]
Since this holds for every bounded measurable \(\varphi\), it is precisely the asserted conditional-mean drift identity \(E_P[\Gamma_q(O;\bar\eta)\mid X=x]-\tau_P(x)=b_q(x)\) for \(P_X\)-almost every \(x\). % lean: hpx_collect % lean: clip_bias

\emph{Step 4: the cancellation mechanisms.}  The closed form
\[
b_q(x)=\{\bar e_q(x)-e_P(x)\}\left\{\frac{\Delta_1(x)}{\bar e_q(x)}+\frac{\Delta_0(x)}{1-\bar e_q(x)}\right\}
\]
is a product of two factors.  If \(\bar e_q(x)=e_P(x)\), the first factor is \(0\); if \(\Delta_0(x)=\Delta_1(x)=0\), the second factor is \(0\).  In either case \(b_q(x)=0\). % lean: clip_bias

\emph{Step 5: no cancellation from \(p_P(x)>q\) alone.}  It remains to show that the stated assumptions do not permit replacing this algebra by a cancellation claim on the region \(\{p_P>q\}\).  Take the covariate space \(\mathbb R\), and let \(P'\) be the law that puts covariate mass \(1\) at \(x=0\), assigns treatment by a fair coin, and returns \(Y=0\) always; explicitly, \(P'\) is the two-atom law
\[
P'=\tfrac12\,\delta_{(0,1,0)}+\tfrac12\,\delta_{(0,0,0)},
\qquad
P'_X=\delta_0,
\qquad
e_{P'}\equiv\tfrac12,
\qquad
\mu_0\equiv\mu_1\equiv 0,
\qquad
\tau_{P'}\equiv 0 .
\]
This \(P'\) is a well-formed observed law with bounded outcomes and satisfies positivity, since \(0<1/2<1\). % lean: clipBiasCounterLaw_wf % lean: clipBiasCounterLaw_bounded % lean: clipBiasCounterLaw_pos
Take \(q'=1/4\) and the plug-in triple \(\bar\mu_0\equiv0\), \(\bar\mu_1\equiv1\), \(\bar e\equiv0\), and evaluate at \(x=0\).  The overlap score is
\[
p_{P'}(0)=\min\{e_{P'}(0),1-e_{P'}(0)\}=\tfrac12>\tfrac14=q',
\]
so \(x=0\) lies strictly inside the region where clipping is claimed not to matter. % lean: clipBiasCounterLaw_overlap
Yet the clipped plug-in propensity is \(\bar e_{q'}(0)=\min\{3/4,\max\{1/4,0\}\}=1/4\), while \(\Delta_1(0)=1-0=1\) and \(\Delta_0(0)=0-0=0\), so the drift formula gives
\[
b_{q'}(0)
=
\left(\tfrac14-\tfrac12\right)
\left\{\frac{1}{1/4}+\frac{0}{3/4}\right\}
=
-\tfrac14\cdot 4
=
-1
\ne 0 .
\]
% lean: clipBiasCounterLaw_bias_ne
Hence \(p_P(x)>q\) alone does not force the drift to vanish, and no additive clipped-region cancellation bound follows from the stated boundedness, positivity, and observed-law conditions alone. % lean: clip_bias
\end{proof}
